% TOP_PROBLEM: {"id":30007608,"problem_number":"LOCAL-30007608","title":"Liquidity backstops for nonbanks","category_id":21,"category":{"id":21,"name":"economics","display_name":"Economics","description":"Open economic theory statements, published research agendas, and proposed scoped specifications; question provenance and historical priority are qualified per record.","slug":"economics","order_index":21,"created_at":"2026-10-08T00:00:00Z"},"status":"research_question","external_url":null,"legacy_ids":[],"published":false,"statement":"In the explicitly specified setting: Which nonbanks should receive central-bank liquidity, on what terms, without undermining private self-insurance? The mathematical statement above fixes the target, assumptions and scope of this catalogue formulation.","economics":{"original_id":97,"field":"Banking, credit and financial stability","kind":"design","formalization_basis":"adapted_research_question","source_status":"explicit_open","priority_status":"earliest_found","priority_note":"This is the earliest explicit statement verified in this audit; first priority for the full catalogue formulation is not established.","earliest_verified_reference":{"authors":["Sirio Aramonte","Andreas Schrimpf","Hyun Song Shin"],"title":"Non-bank financial intermediaries and financial stability","year":2021,"url":"https://www.bis.org/publications/working-paper-972-non-bank-financial-intermediaries-and-financial-stability.pdf","doi":null,"locator":"Section 5, Open questions for policy and research, printed pp.38–39 (PDF pp.41–42), January 2022 revision","evidence_summary":"Explicitly asks optimal oversight and public-backstop design."},"current_reference":{"authors":["Bank of England"],"title":"Bank of England Agenda for Research, 2025–2028","year":2026,"url":"https://www.bankofengland.co.uk/research/bank-of-england-agenda-for-research","doi":null,"locator":"Central Bank Balance Sheet / Central banking toolkit, paragraphs 3–4","evidence_summary":"Requests self-insurance-compatible systemic-liquidity backstop design."},"formalization_note":"Adapts an explicitly verified research-agenda passage. Mathematical objects, interventions, objectives and scope are proposed here; existing model-specific solutions are not relabeled unsolved."},"statusQualification":"Published question or research agenda verified at the cited locator. The mathematical specification is adapted for this catalogue and includes additional assumptions; source publication does not establish first-ever priority or certify this exact model remains unsolved. Adapts an explicitly verified research-agenda passage. Mathematical objects, interventions, objectives and scope are proposed here; existing model-specific solutions are not relabeled unsolved.","statusReviewedAt":"2026-10-08"}
% FIRST_POSTED: 2026-10-08
% ENTRY_KIND: statement_only
% !TeX program = lualatex
\documentclass[11pt]{article}
\usepackage{fontspec}
\usepackage[a4paper,margin=25mm]{geometry}
\usepackage{amsmath,amssymb,mathtools}
\usepackage{xurl}
\usepackage[hidelinks]{hyperref}
\newcommand{\sref}[2]{\hyperlink{source:#1:#2}{[S#2]}}
\setlength{\emergencystretch}{3em}
\begin{document}
% problemId:30007608
\section*{Liquidity backstops for nonbanks}\label{30007608}
\subsection{Definitions and mathematical statement}
Fix nonbank institutions with assets, short-term funding, private liquidity buffers and potential access to central-bank lending. A facility \(a\) specifies eligible institutions, collateral, haircuts, interest charges, borrowing caps and activation conditions. Institution \(i\) chooses a pre-crisis buffer \(b_i(a)\) anticipating the facility. Let \(L(a,\omega;b(a))\) be crisis resource losses after this endogenous response, \(F(a,\omega)\) expected public losses, and \(C(a)\) routine operating costs. The task is to characterize
\[a^*\in\arg\min_{a\in A}E[L(a,\omega;b(a))+F(a,\omega)]+C(a).\]
Here \(A\) is the feasible facility set. Define shock state \(\omega\), its distribution, the institution perimeter and fiscal-risk accounting explicitly. Impose incentive and eligibility constraints ensuring institutions cannot profit from manufacturing collateral or misreporting solvency. Quantify the reduction in liquidity externalities and the offsetting loss of private self-insurance. Compare standing facilities with crisis-only lending while accounting for anticipated discretionary rescues. Determine when access must be coupled to supervision or prudential requirements. State which objects are identified from existing interventions and which require structural assumptions. A liquidity facility that prevents one historical fire sale does not establish an optimal access perimeter or demonstrate that moral hazard is negligible.

\par\textbf{Formulation and scope.} Adapts an explicitly verified research-agenda passage. Mathematical objects, interventions, objectives and scope are proposed here; existing model-specific solutions are not relabeled unsolved.

\par\textbf{Open-question provenance.} An explicit question or research agenda was verified at the source locator. This does not by itself certify that every later version remains unresolved \sref{30007608}{1}.

\par\textbf{Historical priority.} This is the earliest explicit statement verified in this audit; first priority for the full catalogue formulation is not established.

\subsection{Short English statement}
In the explicitly specified setting: Which nonbanks should receive central-bank liquidity, on what terms, without undermining private self-insurance? The mathematical statement above fixes the target, assumptions and scope of this catalogue formulation.

\subsection{Sources}
\begin{itemize}\raggedright
\item[\textnormal{[S1]}]\hypertarget{source:30007608:1}{} Sirio Aramonte, Andreas Schrimpf, Hyun Song Shin. 2021. Non-bank financial intermediaries and financial stability. Section 5, Open questions for policy and research, printed pp.38–39 (PDF pp.41–42), January 2022 revision.
\url{https://www.bis.org/publications/working-paper-972-non-bank-financial-intermediaries-and-financial-stability.pdf}.
{\small\textit{Earliest verified question/agenda reference; historical priority qualified below. Explicitly asks optimal oversight and public-backstop design.}}

\item[\textnormal{[S2]}]\hypertarget{source:30007608:2}{} Bank of England. 2026. Bank of England Agenda for Research, 2025–2028. Central Bank Balance Sheet / Central banking toolkit, paragraphs 3–4.
\url{https://www.bankofengland.co.uk/research/bank-of-england-agenda-for-research}.
{\small\textit{Supporting or current literature; see evidence qualification. Requests self-insurance-compatible systemic-liquidity backstop design.}}
\end{itemize}
\end{document}
